\section{Related Work}\label{sec:related}

\para{The Superficial Alignment Hypothesis.} A growing line of work asks how
much instruction tuning really changes a base model. \citet{zhou2023lima}
articulate the Superficial Alignment Hypothesis: a model's knowledge and
capabilities are learned almost entirely during pretraining, and alignment
mainly teaches it which subdistribution of formats and styles to surface. This
view is supported by the strong performance of LIMA from a thousand curated
examples, and by in-context alignment results suggesting that a few
demonstrations recover much of the aligned behavior \citep{zhao2024iclsufficient,
lin2023urial}. Other work pushes back: \citet{revisiting2024superficial} report
cases where fine-tuning injects new behavior that few-shot prompting cannot
reproduce, arguing the shift is not purely superficial. We do not take a side in
this debate at the level of aggregate behavior. Instead, we operationalize
``superficial'' \emph{per token}: if alignment were entirely a stylistic
re-weighting, the \klt{} at each position would be predictable from cheap
base-side context. Our \gbt{} predictor measures exactly how much of the shift is
this superficial component, and our \emph{residual} isolates the part that is
not.

\para{Token-distribution-shift analysis.} Closest to us is URIAL
\citep{lin2023urial}, which compares aligned and base models token by token. They
introduce a base-rank diagnostic, ranking each token the aligned model emits
within the base model's distribution, and find that the two models already agree
on the top-1 token roughly $78\%$ of the time, that disagreement concentrates in
a small set of stylistic and discourse markers (transition words, safety
disclaimers, formatting), and that the shift decays sharply as a response
proceeds. We reproduce all three findings at the distributional level: top-1
agreement of $0.83$--$0.85$, mean per-token \klt{} of $0.21$--$0.24$ nats, and KL
that decays with position (on \qwensmall{}, mean KL falls from $0.68$ at
positions 0--4 to $0.12$ past position 40). \citet{lin2023urial} read off these
shifted tokens by eye to argue alignment is superficial. Unlike URIAL, we do not
stop at describing \emph{which} tokens shift on average; we fit an explicit
predictor of the shift and ask what is \emph{left over} once the predictable,
stylistic component is removed.

\para{Logit-difference behavior signals.} A parallel thread treats the
difference between aligned and base \emph{logits} as a usable signal rather than
a nuisance. Emulated Fine-Tuning \citep{mitchell2023eft} formalizes the behavior
delta $\log \pi_{\text{inst}} - \log \pi_{\text{base}}$ and shows it can be
transplanted across model scales; proxy-tuning \citep{liu2024proxy} adds a small
tuned model's logit difference to a large base model to steer it without training
the large model. Contrastive-decoding methods exploit the same gap for
generation quality and reasoning: \citet{li2023contrastivedecoding} and
\citet{obrien2023contrastive} contrast a strong and a weak model, and DoLa
\citep{chuang2023dola} contrasts early and late layers of a single model. These
methods \emph{use} the logit difference to decode; they do not ask which parts of
it are anticipatable. Unlike them, we treat the per-token divergence as a
\emph{target to predict} from base-side context, so that the prediction error,
not the difference itself, becomes the object of study.

\para{Probability concentration under alignment.} \citet{yang2026branching}
study alignment through the branching factor and argue that instruction tuning
sharpens the next-token distribution, concentrating mass and surfacing
low-entropy continuations that already existed latently in the base model. This
gives a mechanistic reason why so much of the shift should be predictable from
base-side uncertainty: where the base model is already confident, alignment has
little room to move, and where it is uncertain, alignment often collapses onto a
preferred continuation. Our feature-importance results bear this out, with base
top-1 log-probability and base entropy dominating the predictor. But sharpening
is a statement about the \emph{average} or \emph{typical} token; it does not say
where alignment overrides the base model's preference in ways uncertainty cannot
forecast. That override is precisely what our residual captures.

\para{Transferable weight deltas and interpretability.} Finally, recent work
shows the alignment weight delta is itself a portable object: \citet{shadowft2025}
transfer a fine-tuning delta between models, treating it as a reusable patch.
Interpretability tools such as LogitLens4LLMs \citep{logitlens2025} project
internal states to the vocabulary to read where and how a model commits to a
token. We work in the same spirit of reading the alignment difference directly in
probability space, but at the granularity of individual response positions, and
we use teacher forcing of the \emph{identical} chat-formatted sequence through
both models so that the only thing that varies is the weights.

\para{Positioning.} Across all of this work, the alignment shift is
characterized by its \emph{average} or its \emph{predictable} structure: which
tokens move, how much mass concentrates, how the delta transfers. None of it
trains an explicit predictor of the per-token divergence and then studies the
\emph{residual} of that predictor. That residual, the part of $\klt{}$ that a
learned model genuinely cannot anticipate from simple base-side context, is our
object of study, and we show it is neither noise nor the generic stylistic tokens
everyone has already cataloged, but alignment-specific structure: safety
refusals, discourse and structural pivots, and persona and identity injections.
